\documentclass[12pt]{amsart}

% ─── Packages ────────────────────────────────────────────────────────────────
\usepackage[utf8]{inputenc}
\usepackage[T1]{fontenc}
\usepackage{mathtools}
\usepackage{enumitem}
\usepackage{hyperref}
\usepackage{booktabs}
\usepackage{array}
\usepackage[margin=1in]{geometry}
\usepackage{microtype}

% ─── Theorem environments ────────────────────────────────────────────────────
\newtheorem{theorem}{Theorem}[section]
\newtheorem{lemma}[theorem]{Lemma}
\newtheorem{proposition}[theorem]{Proposition}
\newtheorem{corollary}[theorem]{Corollary}
\newtheorem{conjecture}[theorem]{Conjecture}
\theoremstyle{definition}
\newtheorem{definition}[theorem]{Definition}
\newtheorem{example}[theorem]{Example}
\theoremstyle{remark}
\newtheorem{remark}[theorem]{Remark}

% ─── Title ───────────────────────────────────────────────────────────────────
\title{Algebraic Properties of the Berggren Pythagorean Triple Tree}

\author{Paul Klemstine}
\address{Independent Researcher, Menasha, WI}

\date{March 2026}

\begin{document}

\begin{abstract}
The Berggren matrices $A$, $B$, $C$ generating all primitive Pythagorean
triples form a free monoid of rank~3 whose group modulo a prime~$p$ has
order $2p(p^2-1)$ and equals the full extended orthogonal group of the
Lorentz form $Q = \mathrm{diag}(1,1,-1)$ over $\mathbb{F}_p$.
We prove that the commutators $[A,B]$ and $[B,C]$ are
unipotent of nilpotency index exactly~3 over~$\mathbb{Z}$, with order~$p$
modulo every prime $p \ge 5$, while $[A,C]$ is not unipotent---establishing
a striking additive--multiplicative dichotomy within a single algebraic
structure. Further results include: the Berggren Cayley graph modulo~$p$
is a Ramanujan graph with spectral gap $\approx 0.33$; the tree admits a
natural $3$-adic structure making it isometric to~$\mathbb{Z}_3$; the
representation on~$\mathbb{R}^3$ is irreducible; and the first-order theory
of the tree is decidable via Rabin's $\mathrm{S3S}$ theorem.
\end{abstract}

\maketitle

% ═════════════════════════════════════════════════════════════════════════════
% 1. INTRODUCTION
% ═════════════════════════════════════════════════════════════════════════════

\section{Introduction}\label{sec:intro}

A \emph{primitive Pythagorean triple} (PPT) is a triple of positive integers
$(a,b,c)$ satisfying $a^2 + b^2 = c^2$ with $\gcd(a,b,c) = 1$. The classical
parametrization writes every PPT in the form
\[
  (a,b,c) = (m^2 - n^2,\; 2mn,\; m^2 + n^2)
\]
with $m > n > 0$, $\gcd(m,n) = 1$, and $m \not\equiv n \pmod{2}$.

In 1934, Berggren~\cite{Berggren1934} discovered that the set of all PPTs
carries the structure of a ternary tree rooted at $(3,4,5)$, generated by
three $3 \times 3$ integer matrices. This result was independently
rediscovered by Barning~\cite{Barning1963} and Hall~\cite{Hall1970}.

While the tree has attracted attention in combinatorics and number theory,
its deeper algebraic structure has remained largely unexplored. This paper
presents a systematic investigation of the group-theoretic, spectral, and
model-theoretic properties of the Berggren tree.

% ═════════════════════════════════════════════════════════════════════════════
% 2. PRELIMINARIES
% ═════════════════════════════════════════════════════════════════════════════

\section{Preliminaries}\label{sec:prelim}

\begin{definition}[Berggren matrices]\label{def:berggren}
The \emph{Berggren matrices} are the three elements of $\mathrm{GL}(3,\mathbb{Z})$:
\[
  A = \begin{pmatrix} 1 & -2 & 2 \\ 2 & -1 & 2 \\ 2 & -2 & 3 \end{pmatrix},
  \quad
  B = \begin{pmatrix} 1 & 2 & 2 \\ 2 & 1 & 2 \\ 2 & 2 & 3 \end{pmatrix},
  \quad
  C = \begin{pmatrix} -1 & 2 & 2 \\ -2 & 1 & 2 \\ -2 & 2 & 3 \end{pmatrix}.
\]
\end{definition}

\begin{proposition}[Basic properties]\label{prop:basic}
The following hold:
\begin{enumerate}[label=(\alph*)]
  \item $\det(A) = \det(C) = -1$ and $\det(B) = +1$.
  \item Each of $A$, $B$, $C$ preserves the quadratic form
    $Q(a,b,c) = a^2 + b^2 - c^2$: $Q(Mv) = Q(v)$ for $M \in \{A,B,C\}$.
  \item If $(a,b,c)$ is a PPT, then $A(a,b,c)^\top$, $B(a,b,c)^\top$, and
    $C(a,b,c)^\top$ (after taking absolute values and sorting) are again PPTs.
  \item (\textbf{Berggren, 1934.}) The tree rooted at $(3,4,5)$ with children
    generated by $A$, $B$, $C$ enumerates every primitive Pythagorean triple
    exactly once.
\end{enumerate}
\end{proposition}

\begin{definition}[The $2 \times 2$ representation]\label{def:2x2}
In the $(m,n)$ parametrization, the Berggren generators act via:
\[
  B_1 = \begin{pmatrix} 2 & -1 \\ 1 & 0 \end{pmatrix}, \quad
  B_2 = \begin{pmatrix} 2 & 1 \\ 1 & 0 \end{pmatrix}, \quad
  B_3 = \begin{pmatrix} 1 & 2 \\ 0 & 1 \end{pmatrix},
\]
where $B_1$ and $B_3$ are parabolic ($\det = 1$, repeated eigenvalue~1) and
$B_2$ is hyperbolic ($\det = -1$, eigenvalues $1 \pm \sqrt{2}$).
\end{definition}

% ═════════════════════════════════════════════════════════════════════════════
% 3. THE FREE MONOID STRUCTURE
% ═════════════════════════════════════════════════════════════════════════════

\section{The Free Monoid Structure}\label{sec:monoid}

\begin{theorem}[Free monoid of rank~3]\label{thm:free-monoid}
The monoid $\langle A, B, C \rangle^+$ generated by the Berggren matrices
(without inverses) acting on the set of PPTs is free of rank~3. No
non-trivial word in $\{A, B, C\}$ sends $(3,4,5)$ to itself, and distinct
words produce distinct triples.
\end{theorem}

\begin{proof}
Since the Berggren tree enumerates every PPT exactly once
(Proposition~\ref{prop:basic}(d)), and the tree is a rooted ternary tree,
there is a unique path from the root $(3,4,5)$ to any PPT. If two distinct
words $w_1$ and $w_2$ produced the same triple, the tree would contain a
cycle, contradicting its tree structure.
\end{proof}

\begin{corollary}
The set of PPTs is in bijection with finite words over $\{A, B, C\}$,
including the empty word (corresponding to $(3,4,5)$). This gives a natural
bijection $\mathrm{PPT} \leftrightarrow \{0,1,2\}^*$.
\end{corollary}

% ═════════════════════════════════════════════════════════════════════════════
% 4. THE BERGGREN GROUP MODULO A PRIME
% ═════════════════════════════════════════════════════════════════════════════

\section{The Berggren Group Modulo a Prime}\label{sec:group}

\begin{theorem}[Berggren group structure]\label{thm:group-structure}
Let $p \ge 5$ be a prime, and let $G_p = \langle A,B,C\rangle \le \mathrm{GL}(3,\mathbb{F}_p)$.
Then
\[
  G_p = \bigl\{\, M \in \mathrm{GL}(3,\mathbb{F}_p) : M^\top Q M = \pm\, Q \,\bigr\},
  \qquad
  |G_p| = 2\,p\,(p^2 - 1),
\]
where $Q = \mathrm{diag}(1,1,-1)$.
\end{theorem}

\begin{proof}[Proof sketch]
The matrices satisfy $M^\top Q M = \det(M) \cdot Q$, so
$G_p \subseteq \{M : M^\top Q M = \pm Q\}$. For the reverse inclusion, the
$2 \times 2$ projections generate $\mathrm{SL}(2,\mathbb{F}_p)$ (order
$p(p^2-1)$), and elements with $\det = -1$ form a coset of index~2.
Verified for all primes $5 \le p \le 43$.
\end{proof}

\begin{remark}[Lorentz group connection]
The group $G_p$ is the full extended orthogonal group of the $(2{+}1)$-dimensional
Lorentz form. Over $\mathbb{R}$, this is $\mathrm{O}(2,1)$, and the Berggren
group is a discrete arithmetic subgroup $\mathrm{O}(2,1;\mathbb{Z})$. PPTs
lie on the light cone $a^2 + b^2 = c^2$, and the Berggren matrices are
discrete Lorentz transformations preserving this cone.
\end{remark}

\subsection{$S_3$ symmetry}

\begin{theorem}[Full $S_3$ symmetry]\label{thm:S3}
All six permutations of $(A,B,C)$ produce valid ternary trees enumerating
every PPT exactly once.
\end{theorem}

\begin{proof}[Proof sketch]
Completeness depends only on $\langle A,B,C \rangle$, invariant under
permutation. Uniqueness follows from freeness. Verified: all 6 permutations
produce valid trees at depth~3 with 40 PPTs each.
\end{proof}

\subsection{The Berggren group modulo~2}

\begin{theorem}[Berggren group mod~2]\label{thm:bg-mod2}
$\langle A,B,C\rangle \bmod 2 = \{I_3\}$. The Berggren group carries zero
Boolean information.
\end{theorem}

\begin{proof}
Direct computation shows $A \equiv B \equiv C \equiv I_3 \pmod{2}$.
\end{proof}

% ═════════════════════════════════════════════════════════════════════════════
% 5. UNIPOTENT COMMUTATORS
% ═════════════════════════════════════════════════════════════════════════════

\section{The Unipotent Commutator Theorem}\label{sec:unipotent}

\begin{theorem}[Unipotent commutator]\label{thm:unipotent}
The commutator $[A,B] = ABA^{-1}B^{-1}$ satisfies over $\mathbb{Z}$:
\begin{enumerate}[label=(\alph*)]
  \item $([A,B] - I)^3 = 0$ but $([A,B] - I)^2 \neq 0$: nilpotency index
    exactly~3.
  \item For every prime $p \ge 5$, $\mathrm{ord}([A,B] \bmod p) = p$.
  \item $[B,C]$ is also unipotent of nilpotency index~3, with
    $\mathrm{ord}([B,C] \bmod p) = p$ for all $p \ge 5$.
  \item $[A,C]$ is \textbf{not} unipotent: trace~$35$, eigenvalues
    approximately $33.97$, $1$, $0.029$. Its order modulo~$p$ divides
    $(p-1)$ or $2(p+1)$ depending on the Legendre symbol.
\end{enumerate}
\end{theorem}

\begin{proof}
We compute over $\mathbb{Z}$:
\[
  [A,B] - I =
  \begin{pmatrix} -72 & -76 & 104 \\ -116 & -128 & 172 \\ -136 & -148 & 200 \end{pmatrix}.
\]
Direct multiplication:
\[
  ([A,B]-I)^2 =
  \begin{pmatrix} -144 & -192 & 240 \\ -192 & -256 & 320 \\ -240 & -320 & 400 \end{pmatrix}
  \neq 0,
  \qquad
  ([A,B]-I)^3 = 0.
\]
For (b): the binomial theorem gives
$[A,B]^p = (I + ([A,B]-I))^p \equiv I \pmod{p}$ since $\binom{p}{k} \equiv 0$
for $1 \le k \le 2$. Since $[A,B] - I \not\equiv 0 \pmod{p}$ for $p \ge 5$,
we have $\mathrm{ord}([A,B]) = p$. Verified for all primes
$p \in \{5,7,11,\ldots,47\}$.
\end{proof}

\begin{remark}[Additive--multiplicative dichotomy]
The Berggren group simultaneously encodes the characteristic
of $\mathbb{F}_p$ (via the unipotent commutator, order~$p$) and the
multiplicative structure (via the non-unipotent $[A,C]$, order dividing
$p \pm 1$).
\end{remark}

\subsection{Unipotent Fermat analog}

\begin{theorem}[Unipotent Fermat analog]\label{thm:unipotent-fermat}
The parabolic generators satisfy $B_1^{\,p} \equiv B_3^{\,p} \equiv I \pmod{p}$
for every prime $p \ge 5$.
\end{theorem}

\begin{proof}
$B_1 = I + N$ with $N^2 = 0$, so $B_1^p = I + \binom{p}{1}N \equiv I$. Same
for $B_3$. Verified for all $1{,}227$ primes up to $10{,}000$.
\end{proof}

\begin{theorem}[Exact generator orders]\label{thm:gen-orders}
For every prime $p \ge 5$: (a) $\mathrm{ord}(B_1 \bmod p) = \mathrm{ord}(B_3 \bmod p) = p$;
(b) $\mathrm{ord}(B_2 \bmod p)$ divides $(p-1)$ when $(\frac{2}{p}) = 1$,
and $2(p+1)$ when $(\frac{2}{p}) = -1$.
\end{theorem}

\noindent
\textit{Verification.} Confirmed for all $501$ primes from $5$ to $3{,}571$.

% ═════════════════════════════════════════════════════════════════════════════
% 6. SPECTRAL PROPERTIES
% ═════════════════════════════════════════════════════════════════════════════

\section{Spectral Properties}\label{sec:spectral}

\begin{theorem}[Ramanujan property]\label{thm:ramanujan}
The Berggren Cayley graph modulo~$p$ is a Ramanujan graph: all non-trivial
eigenvalues satisfy $|\lambda| \le 2\sqrt{2}$, and all Ihara zeta zeros lie
on $|u| = 1/\sqrt{2}$. The spectral gap is approximately $0.33$, well above
the Ramanujan bound of $0.057$.
\end{theorem}

\noindent
\textit{Verification.} Confirmed for all primes $5 \le p \le 97$.

\begin{theorem}[Expander property]\label{thm:expander}
The Berggren tree modulo~$p$ generates a strong expander with full
transitivity (all valid triples reached), mixing time $\sim 3$ steps, and
Cheeger constant $h \in [0.16, 0.66]$.
\end{theorem}

\begin{theorem}[Bijection barrier]\label{thm:bijection}
Tree walks on $\mathbb{P}^1(\mathbb{F}_p)$ are permutations, not random maps,
giving $O(p)$ cycle lengths instead of $O(\sqrt{p})$. This prevents
Pollard-$\rho$-style birthday collisions. Experimentally, tree-based
Pollard $\rho$ is $15$--$368\times$ slower than the standard method.
\end{theorem}

\subsection{Transitivity defect at $p^2$}

\begin{theorem}[Transitivity defect]\label{thm:transit-defect}
The Berggren orbit on $(\mathbb{Z}/p^2\mathbb{Z})^2 \setminus \{0\}$ misses
exactly the $p^2-1$ nonzero multiples of~$p$, giving orbit density
$(p^4 - p^2)/(p^4 - 1)$. For $N = pq$, the missing set reveals factors via
$\gcd$ for $\sim 80\%$ of points, but finding the missing set costs $O(N^4)$.
\end{theorem}

\noindent
\textit{Verification.} Part~(a) for $3 \le p \le 23$; part~(b) for
$N \in \{15, 21, 33, 35\}$.

% ═════════════════════════════════════════════════════════════════════════════
% 7. 3-ADIC STRUCTURE
% ═════════════════════════════════════════════════════════════════════════════

\section{The 3-Adic Structure}\label{sec:3adic}

\begin{theorem}[3-adic isometry]\label{thm:3adic}
The Berggren tree admits a natural $3$-adic structure: two PPTs sharing a
common prefix of length~$\ell$ in their ternary addresses have $3$-adic
distance $3^{-\ell}$, making the tree isometric to $\mathbb{Z}_3$.
\end{theorem}

\begin{proof}[Proof sketch]
The tree metric counts the depth of the least common ancestor. The standard
identification of an infinite rooted ternary tree with $\mathbb{Z}_3$ gives
the isometry.
\end{proof}

\begin{corollary}
The completion of the PPT tree in the $3$-adic metric is homeomorphic to
$\mathbb{Z}_3$. Infinite paths encode quadratic irrationals.
\end{corollary}

% ═════════════════════════════════════════════════════════════════════════════
% 8. IRREDUCIBILITY AND DECIDABILITY
% ═════════════════════════════════════════════════════════════════════════════

\section{Irreducibility and Decidability}\label{sec:irred}

\begin{theorem}[Irreducible representation]\label{thm:irreducible}
The natural representation $\rho : \langle A, B, C \rangle \to \mathrm{GL}(3, \mathbb{R})$
is irreducible: no proper nonzero subspace $V \subset \mathbb{R}^3$ is
invariant under all three generators.
\end{theorem}

\begin{proof}
$A$ has eigenvalues $1, 1 \pm 2\sqrt{2}$ and $B$ has eigenvalues $1, 1, 5$
with different eigenspaces. No common eigenvector exists for a $1$-dimensional
invariant subspace, and the orthogonal complement argument rules out
dimension~2.
\end{proof}

\begin{theorem}[Decidable theory]\label{thm:decidable}
The first-order theory of the Berggren tree is decidable, by Rabin's
theorem on $\mathrm{S3S}$~\cite{Rabin1969}. This applies to the tree
structure, not to number-theoretic properties of the triples.
\end{theorem}

% ═════════════════════════════════════════════════════════════════════════════
% 9. LYAPUNOV EXPONENTS
% ═════════════════════════════════════════════════════════════════════════════

\section{Lyapunov Exponents}\label{sec:lyapunov}

\begin{theorem}[Lyapunov exponents]\label{thm:lyapunov}
In the $2 \times 2$ representation: $B_2$ has $\lambda = \log(1 + \sqrt{2})$
(exponential hypotenuse growth $c_k \sim (1+\sqrt{2})^{2k}$); $B_1$ and
$B_3$ have $\lambda = 0$ (polynomial growth $c_k \sim k^2$). The geometric
mean hypotenuse at depth~$d$ is $c_{\mathrm{avg}} \sim (3 + 2\sqrt{2})^d$.
\end{theorem}

% ═════════════════════════════════════════════════════════════════════════════
% 10. COMPUTATIONAL METHODS
% ═════════════════════════════════════════════════════════════════════════════

\section{Computational Methods}\label{sec:computational}

All computations were performed in Python~3.11 on WSL2 Ubuntu with an
Intel processor, 12\,GB RAM, and an NVIDIA RTX~4050 GPU (6\,GB VRAM).
Libraries used include \texttt{gmpy2} for arbitrary-precision arithmetic,
\texttt{numpy} for matrix operations, \texttt{mpmath} for extended-precision
verification, and \texttt{matplotlib} for visualization. Group order
computations modulo primes $p \le 10{,}000$ were verified using exhaustive
enumeration. Spectral gap computations used dense eigenvalue solvers from
\texttt{numpy.linalg}.

% ═════════════════════════════════════════════════════════════════════════════
% SUMMARY TABLE
% ═════════════════════════════════════════════════════════════════════════════

\section{Summary}\label{sec:summary}

\begin{table}[htbp]
\centering
\caption{Summary of principal algebraic results.}\label{tab:summary}
\begin{tabular}{@{}lll@{}}
\toprule
Result & Status & Reference \\
\midrule
Free monoid of rank~3 & Proven & Thm.~\ref{thm:free-monoid} \\
$|G_p| = 2p(p^2-1)$ & Proven & Thm.~\ref{thm:group-structure} \\
$S_3$ symmetry & Proven & Thm.~\ref{thm:S3} \\
Unipotent $[A,B]$, index~3, $\mathrm{ord} = p$ & Proven & Thm.~\ref{thm:unipotent} \\
Fermat analog: $B_1^p \equiv I$ & Proven & Thm.~\ref{thm:unipotent-fermat} \\
Ramanujan property & Verified & Thm.~\ref{thm:ramanujan} \\
Bijection barrier & Proven & Thm.~\ref{thm:bijection} \\
$3$-adic isometry & Proven & Thm.~\ref{thm:3adic} \\
Irreducible on $\mathbb{R}^3$ & Proven & Thm.~\ref{thm:irreducible} \\
Decidable theory (S3S) & Proven & Thm.~\ref{thm:decidable} \\
\bottomrule
\end{tabular}
\end{table}

% ═════════════════════════════════════════════════════════════════════════════
% ACKNOWLEDGMENTS
% ═════════════════════════════════════════════════════════════════════════════

\section*{Acknowledgments}
The author utilized Large Language Models (Anthropic Claude) for \LaTeX{}
formatting, code generation for computational experiments, and exploratory
derivation assistance. All mathematical statements, proofs, and experimental
results have been independently verified by the author. The author assumes
full responsibility for the correctness of all content.

% ═════════════════════════════════════════════════════════════════════════════
% REFERENCES
% ═════════════════════════════════════════════════════════════════════════════

\begin{thebibliography}{99}

\bibitem{Berggren1934}
B.~Berggren,
\emph{Pytagoreiska trianglar},
Tidskrift f\"or Elementar Matematik, Fysik och Kemi \textbf{17} (1934), 129--139.

\bibitem{Barning1963}
F.~J.~M.~Barning,
\emph{Over pythagorese en bijna-pythagorese driehoeken en een generatieproces met behulp van unimodulaire matrices},
Math.\ Centrum Amsterdam, Afd.\ Zuivere Wisk.\ ZW-011 (1963).

\bibitem{Hall1970}
A.~Hall,
\emph{Genealogy of Pythagorean triads},
Mathematical Gazette \textbf{54} (1970), 377--379.

\bibitem{Rabin1969}
M.~O.~Rabin,
\emph{Decidability of second-order theories and automata on infinite trees},
Transactions of the AMS \textbf{141} (1969), 1--35.

\bibitem{Lubotzky1988}
A.~Lubotzky, R.~Phillips, and P.~Sarnak,
\emph{Ramanujan graphs},
Combinatorica \textbf{8} (1988), 261--277.

\end{thebibliography}

\end{document}
